\section{The problem and the honest position}

Write $p_n$ for the $n$-th prime and $H_1=\liminf_{n\to\infty}(p_{n+1}-p_n)$.
The twin prime conjecture is the statement $H_1=2$. It is \lbl{Open}.

A finite set $\mathcal H=\{h_1<\dots<h_k\}$ of integers is \emph{admissible} if
for every prime $p$ it misses at least one residue class modulo $p$. Let
$\DHL[k,2]$ be the statement that every admissible $k$-tuple has infinitely
many translates $n+\mathcal H$ containing at least two primes, and let $H(k)$
be the least diameter $h_k-h_1$ of an admissible $k$-tuple. Then
\begin{equation}\label{eq:reduction}
\DHL[k,2]\ \Longrightarrow\ H_1\le H(k),
\end{equation}
and $\DHL[2,2]$ applied to $\{0,2\}$ is the twin prime conjecture. All progress
on $H_1$ since 2013 consists of proving $\DHL[k,2]$ for smaller $k$.

\subsection*{Where the bound stands (13 November 2026)}

\begin{center}
\begin{tabular}{@{}rlll@{}}
\toprule
$H_1\le$ & $k$ & Source & Label\\
\midrule
$7\times10^7$ & $3.5\times10^6$ & Zhang \cite{Zha14} & \lbl{Literature}\\
600 & 105 & Maynard \cite{May15} & \lbl{Literature}\\
246 & 50 & Polymath8b \cite{Pol14b} & \lbl{Literature}\\
240 & 49$^{\ast}$ & Stadlmann \cite{Sta26} & \lbl{Claimed}$^{\dagger}$\\
236 & 48 & Kintali \cite{Kin26}; Song--Yue \cite{SY26} & \lbl{Claimed}$^{\ddagger}$\\
212 & 45 & Charton et al.\ \cite{Cha26} & \lbl{Claimed}\\
186 & 40 & OpenAI \cite{OAI26} & \lbl{Claimed}\\
182 & 39 & GPT-6 Astra, prompted by Sharifi \cite{Sha26} & \lbl{Claimed}\\
\midrule
12 & 5 & Maynard, assuming EH \cite{May15} & \lbl{Literature}\\
6 & 3 & Polymath8b, assuming GEH \cite{Pol14b} & \lbl{Literature}\\
\bottomrule
\end{tabular}
\end{center}
\noindent{\small $^{\ast}$ Value of $k$ inferred by us from $H(49)=240$ and $H(48)=236$, not read from the source.
$^{\dagger}$ Its variational step is reproduced here (Corollary~\ref{cor:240}); its equidistribution input is not verified by us. Song and Yue \cite[Remark~4.23]{SY26} point out a possible missing hypothesis in the proof of \cite[Lemma~5]{Sta26} (a Type~IIa lemma applied after swapping the factors); they say it is not a counterexample to the conclusion. In version~0.10 we read the cited step in the arXiv text (v1): in the range $\gamma\in(\frac12,\frac12+2\omega+\varepsilon]$ the proof of \cite[Lemma~5]{Sta26} applies \cite[Lemma~3]{Sta26} with the factors swapped, Lemma~3 needs its shorter factor to be Siegel--Walfisz, and neither Lemma~5 nor case~(I) of \cite[Definition~9]{Sta26} assumes this of $\alpha$. So the step reads as they say. Whether the coefficient sequences that actually occur satisfy the missing hypothesis, which would repair the step, we have not checked (Remark~\ref{rem:sta-lemma5}). Both $236$ preprints use the same support (``D182''): $k=48$ is read from \cite{Kin26,SY26}.
$^{\ddagger}$ Its variational step is reproduced here (Corollary~\ref{cor:236}); its equidistribution input is not verified by us.}

The bound $246$ stood from 2014 until the last days of August 2026. It has been
formalised in Lean~4, assuming only the Bombieri--Vinogradov theorem
\cite{Axiom26}. The bounds from $240$ to $186$ appeared within one week.
They rest on Stadlmann's idea of enlarging the support of the sieve weights
using equidistribution estimates for moduli with a large smooth factor. The
Lean development accompanying the bound $186$ is conditional on three stated
inputs (two exponential-sum bounds and a family of numerical integral bounds)
\cite{OAI26}. An independent re-run of that repository, dated 7 and 8 November
2026, reproduced the build and the numerical certificate and describes the
result as ``conditionally supported, with the computation independently
reproduced; not yet settled''; the two exponential-sum bounds remain axioms
\cite{NM26}. A further document, dated 18 September 2026 and naming an AI
system as its author, claims $\DHL[39,2]$ and hence $182$; by its own abstract
its Lean development verifies the last implication ``under explicit numerical
and finite-field premises'' \cite{Sha26}. We had missed it in earlier
versions. (These details were read by an earlier run of this project; on
12 November the file was too large for our reader and direct download was
refused, so we have not re-checked them since.) As of 7 September 2026 a standard reference still listed $246$ as
the best fully established value \cite{MW}, and the page read the same on 12
November. On the date of this version we found no report of complete
independent verification of any of the 2026 claims and no claimed bound below
$182$. We therefore keep three numbers apart throughout: \emph{established}
($246$), \emph{claimed} ($182$), and \emph{reproduced by this project}:
$246$ unconditionally; since version~0.6 also $240$, conditionally on
the equidistribution estimates of \cite{Sta26}, which we have not verified
(Corollary~\ref{cor:240}); and since version~0.11 also $236$, conditionally
on the equidistribution theorem of \cite{SY26} (Corollary~\ref{cor:236}),
which we have not verified either.

\subsection*{Why \texorpdfstring{$2$}{2} is a different problem}
Under the generalised Elliott--Halberstam conjecture the sieve gives $H_1\le6$,
and Polymath8b showed that $6$ cannot be improved by purely sieve-theoretic
arguments: this is Selberg's parity obstruction \cite{Pol14b}. Lowering the
bound (Track~A below) and proving the conjecture (Track~B) are therefore
different tasks. No amount of success on Track~A reaches $2$.

